% paper/sections/08_time_integration.tex
\section{時間積分・移流スキームの高精度化}
\label{sec:time}

\subsection{WENO5 による移流：なぜ必要か}
\label{sec:weno5}

対流項および CLS の移流項には，5次精度の Weighted ENO (WENO5) スキーム \cite{JiangShu1996} を採用する．

\subsubsection{気液界面移流の困難}

気液界面の移流は，数値スキームの設計において特殊な問題を提起する．
$\psi$ は界面付近で 0 から 1 へと急激に変化するが，この勾配の強さは
物理的実体（物性の不連続）を反映したものであって，数値拡散で
「なだらかにして良い」ものではない．

通常の線形スキーム（例：4次中心差分）は高精度だが，
$\psi$ のような急峻なプロファイルを移流すると
\textbf{Gibbs 振動}（負の $\psi$ や $\psi > 1$ の非物理的な値）が発生する．
これが界面の「にじみ」（1章・失敗例4）の主な原因の一つである．

\subsubsection{ENO・WENO の基本思想}

\textbf{ENO（Essentially Non-Oscillatory）法}の思想は，「局所的な滑らかさを検出し，
不連続（または急峻な勾配）を含まない滑らかな候補ステンシルを選んで使う」ことである．

WENO 法はこれを改良し，複数のステンシルを\textbf{滑らかさ指標（smoothness indicator）
に基づく重みで凸結合}する：
\begin{equation}
  \hat{f}^{\text{WENO}} = \sum_{k=0}^{r-1} \omega_k \, \hat{f}_k
  \label{eq:weno_convex}
\end{equation}
ここで $\hat{f}_k$ は $k$ 番目の候補ステンシルからの数値フラックス，
$\omega_k$ は滑らかさ指標 $\beta_k$ から導出される重みで，
$\sum_k \omega_k = 1$，$\omega_k \ge 0$（凸結合）を満たす．

\textbf{WENO5 の場合（5次精度，3つのステンシル）：}
\begin{itemize}
  \item 3つの3点ステンシル（それぞれ3次精度）を凸結合し，滑らかな領域では5次精度を達成
  \item 不連続付近では，その不連続を含むステンシルの重みが自動的にゼロに近づき，
        滑らかなステンシルのみが使われる（振動抑制）
\end{itemize}

\begin{tcolorbox}[mybox, title={WENO5 の滑らかさ指標（1次元 $x$ 方向，上流フラックス）}]
格子点 $i$ における3つの候補ステンシルの滑らかさ指標 $\beta_k$ は：
\begin{align}
  \beta_0 &= \frac{13}{12}(f_{i-2}-2f_{i-1}+f_i)^2 + \frac{1}{4}(f_{i-2}-4f_{i-1}+3f_i)^2 \notag\\
  \beta_1 &= \frac{13}{12}(f_{i-1}-2f_i+f_{i+1})^2 + \frac{1}{4}(f_{i-1}-f_{i+1})^2 \notag\\
  \beta_2 &= \frac{13}{12}(f_i-2f_{i+1}+f_{i+2})^2 + \frac{1}{4}(3f_i-4f_{i+1}+f_{i+2})^2
  \label{eq:weno5_beta}
\end{align}
$\beta_k$ が小さいほど，そのステンシルの範囲が滑らか（$\Rightarrow$ 高い重み $\omega_k$）．
不連続を含むステンシルでは $\beta_k$ が大きくなり，重みがほぼゼロになる．
\end{tcolorbox}

\begin{tcolorbox}[defbox, title={WENO5 の性質まとめ}]
\begin{itemize}
  \item 滑らかな領域：5次精度（$O(h^5)$）の空間精度
  \item 不連続・急峻な勾配付近：重みにより振動を自動的に抑制（ENO 的挙動）
  \item 保存形：フラックス形式で書けるため，離散的質量保存が成立
  \item 移流方向の上流化：情報伝播方向を考慮した上流側ステンシルを優先
\end{itemize}
\end{tcolorbox}

\subsection{時間積分スキーム}
\label{sec:time_int}

\subsubsection{3次 TVD Runge--Kutta 法（移流・CLS）}

CLS 移流および NS 方程式の対流項には，Shu--Osher 型 TVD-RK3 法 \cite{ShuOsher1988} を用いる：
\begin{align}
  q^{(1)} &= q^n + \Delta t \, \mathcal{L}(q^n) \notag\\
  q^{(2)} &= \frac{3}{4}q^n + \frac{1}{4}q^{(1)} + \frac{1}{4}\Delta t\,\mathcal{L}(q^{(1)}) \notag\\
  q^{n+1} &= \frac{1}{3}q^n + \frac{2}{3}q^{(2)} + \frac{2}{3}\Delta t\,\mathcal{L}(q^{(2)})
  \label{eq:tvd_rk3}
\end{align}
ここで $\mathcal{L}(q)$ は空間演算子（WENO5 フラックスの発散）を表す．

\begin{tcolorbox}[defbox, title={TVD の意味：全変動を増大させない}]
TVD（Total Variation Diminishing）とは，数値解の全変動
$TV(q) \equiv \sum_i |q_{i+1} - q_i|$ が時間的に増大しない性質を指す：
$TV(q^{n+1}) \leq TV(q^n)$．
これにより，$\psi$ のプロファイルに新たな振動が生成されることが防止され，
界面のギザギザな数値ノイズが自然に抑制される．
特に CLS 変数 $\psi \in [0,1]$ の範囲外への逸脱を防ぐ効果がある．
\end{tcolorbox}

\subsubsection{Crank--Nicolson 法（粘性項）}

粘性項（2階空間微分を含む）は，純陽的処理では厳しい時間刻み制限（$\Delta t < O(h^2)$）を課す．
これを緩和するため，Crank--Nicolson（CN）法による\textbf{半陰的}処理を採用する：
\begin{equation}
  \frac{u^{n+1} - u^n}{\Delta t} = \frac{1}{2}\bigl[\mathcal{D}(u^n) + \mathcal{D}(u^{n+1})\bigr]
  \label{eq:crank_nicolson}
\end{equation}
ここで $\mathcal{D}$ は粘性拡散演算子．CN 法は時間方向に2次精度（$O(\Delta t^2)$）かつ
\textbf{無条件安定}（スペクトル半径の制約なし）であり，適度な $\Delta t$ での安定計算が可能となる．

\subsection{再初期化の Godunov 型上流化}
\label{sec:godunov}

CLS 法の再初期化あるいは標準 Level Set の再初期化において，
Eikonal 方程式の求解には Godunov 型の上流差分が必要となる．
情報伝播方向 $s=\operatorname{sgn}(\phi_0)$ に応じた統一表現は以下の通りである：
\begin{equation}
 \norm{\bnabla\phi}^{\mathrm{uw}}_{i,j}=
 \sqrt{G_x^2+G_y^2},\qquad G_x^2=\max(s\,D_x^-,0)^2+\min(s\,D_x^+,0)^2
 \label{eq:godunov_upwind}
\end{equation}
ここで $D^\pm$ は前進・後退差分である．

\subsection{安定性制約（時間刻みの決定）}
\label{sec:stability}

時間刻み $\Delta t$ は，以下の3種類の制約のうち最も厳しい条件を選ぶ（安全率 $\text{CFL}_{\max} < 1$）：

\begin{tcolorbox}[mybox, title={時間刻みの3制約}]
\begin{align}
  &\text{\textbf{対流 CFL}：}\quad
  \Delta t_{\text{adv}} = \text{CFL}_{\max} \cdot \frac{1}{\max\!\left(\dfrac{|u|_{\max}}{\Delta x}+\dfrac{|v|_{\max}}{\Delta y}\right)}
  \label{eq:dt_adv} \\[6pt]
  &\text{\textbf{粘性 CFL}：}\quad
  \Delta t_{\text{visc}} = \frac{\min(\Delta x, \Delta y)^2}{4(\mu/\rho)_{\max}}
  \quad \text{（陽的処理の場合のみ適用）}
  \label{eq:dt_visc} \\[6pt]
  &\text{\textbf{毛管波制約}：}\quad
  \Delta t_\sigma = \sqrt{\frac{(\rho_l + \rho_g)\min(\Delta x, \Delta y)^3}{4\pi\sigma}}
  \label{eq:dt_sigma}
\end{align}
最終的な時間刻み：$\Delta t = \min(\Delta t_{\text{adv}},\, \Delta t_{\text{visc}},\, \Delta t_\sigma)$
\end{tcolorbox}

\begin{tcolorbox}[warnbox, title={毛管波制約 $\Delta t_\sigma \propto \sqrt{\Delta x^3}$ が厳しい理由}]
表面張力が支配的な問題（$We \ll 1$，小液滴，毛管現象）では，
毛管波（capillary wave）の位相速度が大きくなり CFL 条件が厳しくなる．
特に $\Delta t_\sigma \propto \Delta x^{3/2}$ の依存性（格子を半分にすると $\approx 2.8$ 倍制約が厳しくなる）
のため，高解像度計算では毛管波制約が支配的になることが多い．

\textbf{対策：}
\begin{itemize}
  \item 表面張力の半陰的処理（Implicit Surface Tension）による制約の緩和
  \item 問題によっては圧力 $p$ から $p + p_0$（静水圧 $p_0 = \rho g z$）を分離する
        動的圧力（kinematic pressure）分割で寄生流れを抑制しつつ $\Delta t$ を拡大
\end{itemize}
\end{tcolorbox}
